\subsection{Multiplying Polynomials: Rectangle Method} 
We can think of multiplication as describing the \makeblank[1in] of a rectangle.
\begin{Example}Multiply $(3x^2+5x+4)(2x^2+6x+3)$ by interpreting the polynomials as describing the sides of a rectangle.
\begin{lpic}{}{6}
\draw (2,-1) rectangle ++(9,-3);
%\foreach \x in {3,6} \draw (2+\x,-1) -- ++(0,-3);
%\foreach \y in {1,2} \draw (2,-1-\y) -- ++(9,0);
\foreach \x/\lbl in {0/3x^2,1/+5x,2/+4} \regtext{2+1.5+3*\x}{-1}{$\lbl$};
\foreach \y/\lbl in {0/2x^2,1/+6x,2/+3} \regtextright{2}{-2-\y}{$\lbl$};
\end{lpic}
\end{Example}
\begin{Example}Multiply $(4x+3y-6)(3x-5y+2)$.
\begin{lpic}{}{6}
\end{lpic}
\end{Example}